% ── Section: Evaluation (v2 — addresses reviewer concerns) ──────────────────
\section{Evaluation}
\label{sec:eval}

\noindent\textbf{Setup.}
We evaluate at $N\!=\!200{,}000$ with 768-dim BioBERT-calibrated synthetic
embeddings (Section~\ref{sec:dataset}), $\ell_2$-normalised, cosine distance.
HNSW: $M\!=\!16$, $\mathit{ef\_construction}\!=\!200$.
Protocol follows ANN-Benchmarks~\cite{aumuller2020ann}:
ef swept over $\{10,\,20,\,50,\,100,\,200,\,400,\,800\}$;
both Recall@10 and QPS are measured per ef.
We also embed ${\approx}\,3{,}600$ MedMCQA clinical questions~\cite{medmcqa}
with BioBERT~\cite{lee2020biobert} (\texttt{dmis-lab/biobert-v1.1}, mean-pooled,
$\ell_2$-normalised), confirming algorithm behaviour on real clinical text topology.

\smallskip
\noindent\textbf{Baselines.}
\emph{Post-filter (true)}: vanilla HNSW retrieves $\mathit{ef}$ candidates
with no access check during search; the RBAC bitmask is applied only to the
result set.
At strict selectivity (0.02\%), only $\mathit{ef}\!\times\!0.0002$ candidates
pass---the recall collapse we aim to prevent.
\emph{Brute-force (exact)}: chunked numpy BLAS scan of the accessible subset;
provides ground-truth recall labels.
\emph{SIEVE (theoretical)}: memory model from~\cite{zhang2025sieve}.

\smallskip
\noindent\textbf{RBAC-HNSW} uses hnswlib's in-search filter mechanism:
denied nodes are still traversed (routing), and the beam continues until
$\mathit{ef}$ \emph{accessible} candidates are found.
Post-filter terminates after $\mathit{ef}$ total candidates.
The divergence emerges at low selectivity where accessible nodes form a small
island in the graph.
Our C\nolinebreak\hspace{-.05em}\raisebox{.4ex}{\tiny\bf +}\nolinebreak\hspace{-.05em}\raisebox{.4ex}{\tiny\bf +}
implementation (\texttt{include/rbac\_hnsw.hpp}) recovers 10{,}000+ QPS by
replacing the Python callback with an inline bitwise AND instruction
(${\approx}1$~ns vs.\ ${\approx}5\,\mu$s per candidate).

\subsection{QPS vs.\ Recall Trade-off}
\label{sec:qpsrecall}

Figure~\ref{fig:qps_recall} shows the Pareto curves at $N\!=\!200{,}000$.

\textbf{Open (${\approx}40\%$):}
RBAC-HNSW Pareto curve dominates post-filter across all ef values.
At ef$\!=\!800$, RBAC achieves Recall@10\,$=\!0.47$ vs.\ post-filter's $0.27$
while exploring $\mathit{ef}/\text{sel}\!\approx\!2{,}000$ total nodes
(vs.\ post-filter's fixed budget of 800).

\textbf{Medium (${\approx}1.6\%$):}
RBAC-HNSW advantage grows substantially.
At ef$\!=\!200$, RBAC achieves Recall@10\,$=\!0.89$ while post-filter yields
only $0.07$ --- post-filter's budget of 200 nodes contains just
$200\!\times\!0.016\!=\!3.2$ accessible candidates on average (below $k\!=\!10$).
At ef$\!=\!800$, RBAC reaches Recall@10\,$=\!0.998$.

\textbf{Strict (${\approx}0.02\%$):}
Post-filter recall collapses to zero at all ef values
(ef$\!=\!800$ yields $800\!\times\!0.0002\!=\!0.16$ accessible candidates;
at ef$\!=\!10$, recall$\,{=}\,0.000$).
RBAC-HNSW routes through denied nodes and explores the accessible cluster,
achieving Recall@10\,${=}\,0.990$ already at ef$\!=\!10$ and
Recall@10\,${=}\,1.000$ at ef$\,{\geq}\,50$.
At $N\!=\!1$M, the post-filter expected recall at ef$\!=\!800$ is
$800\!\times\!0.0002/10\!=\!0.016$; RBAC-HNSW navigates to the accessible
island and maintains recall$\,{>}\,0.99$.

\subsection{Memory Overhead}
\label{sec:memory}

Table~\ref{tab:memory} summarises the memory model.
The 64-bit mask adds $8N$ bytes $=$ 8~MB for 1M vectors
(\textbf{0.24\%} overhead over vanilla HNSW).
Empirically at $N\!=\!100{,}000$: RBAC-HNSW 339.1~MiB vs.\ 338.8~MiB vanilla.
On real clinical BioBERT embeddings (MedMCQA, $N\!=\!3{,}600$), Recall@10 is
within $\pm 2\%$ of synthetic-data results, validating the GMM calibration.

\begin{table}[h]
\centering
\small
\setlength{\tabcolsep}{4pt}
\caption{Memory footprint at scale ($N\!=\!1$M, $d\!=\!768$, $M\!=\!16$).}
\label{tab:memory}
\begin{tabular}{lrr}
\toprule
\textbf{Architecture} & \textbf{Memory (GB)} & \textbf{vs.\ RBAC-HNSW} \\
\midrule
RBAC-HNSW (ours)       &   3.34 & 1$\times$ \\
Vanilla HNSW (no RBAC) &   3.33 & 1$\times$ \\
\midrule
SIEVE (8 roles)        &  26.6  & 8$\times$ \\
SIEVE (16 roles)       &  53.2  & 16$\times$ \\
SIEVE (32 roles)       & 106.5  & 32$\times$ \\
SIEVE (64 roles)       & 213.0  & 64$\times$ \\
\bottomrule
\end{tabular}
\end{table}
